\cotefile{157-4}
\scanpage{98}
\mspage{91}

Soit

\[
 C=\deleted{\unsure{\mathcal M}}\unsure{\underline{M}|\underline{\mathcal M}}^{\pi}/\unsure{\mathscr X|\mathfrak X}
\]

la catégorie suivante. Les objets sont les triples

\[
 (Z,i,\unsure{u|\alpha})\mid Z\in M^{\pi},\quad i\in I,\quad \alpha:Z\longrightarrow X_i,
\]

les flèches

\[
 (Z,i,\alpha)\longrightarrow(Z',i',\alpha')
\]

sont les couples $(u,\lambda)$ avec

\[
 u:Z\longrightarrow Z',\qquad \lambda:i\longrightarrow i'
\]

(i.e. on a $i\leq i'$, et $\lambda$ est l'unique flèche de $i$ dans $i'$) tels que le diagramme

\diagram{\begin{tikzcd}[ampersand replacement=\&]
 Z\arrow[r,"\alpha"]\arrow[d,"u"'] \& X_i\arrow[d,"{\text{transition }\lambda_{i',i}}"]\\
 Z'\arrow[r,"\alpha'"'] \& X_{i'}
\end{tikzcd}}

commute. On \unsure{peut} donc \unsure{définir} $C$ comme le produit fibré dans le carré cartésien

\diagram{\begin{tikzcd}[ampersand replacement=\&,column sep=large]
 {} \& C\arrow[dl,dashed,"{\unsure{\mathscr Z|Z}}"']\arrow[d]\arrow[r,"\text{cofibré}"] \& I\arrow[d,"{\mathscr X:\unsure{X|i}\longmapsto X_i}"]\\
 \underline{M}^{\pi} \& \underline{M}^{\pi}/E\arrow[l,dashed,"\text{source}"]\arrow[r,"\text{but}","\text{cofibré}"'] \& E
\end{tikzcd}}

\deleted{\unsure{La|comme}} La fibre \deleted{\unsure{d'}} de $C$ en $i\in\deleted{\unsure{E}} I$ est $\unsure{\mathrm{Filt}|\mathrm{FétUt}|\mathrm{FibUt}}^{\pi}/X_i$. On a \deleted{\unsure{comme}} \inserted{(comme b) $\Rightarrow$ \deleted{\unsure{que} \unsure{problème|hypothèse} \unsure{pr} \unsure{des}} \unsure{première|premier|promotion} \unsure{moitié|critère|critéri} des \unsure{$\varphi$},)} que $\underline M^{\pi}$ \deleted{\unsure{en}} est \unsure{un|une}

% archive sheet number: 97
